\documentclass{article}

\usepackage{../packages}
\usepackage{../macros}

\title{Virtual Sets}
\author{mark}

\begin{document}
\maketitle

% \section{Grothendieck group for Monoidal Categories}

% \begin{definition}
%     A rig category is a category $\cC$ equipped with 
%     \begin{enumerate}
%         \item Two symmetric monoidal structures $(+, 0, a, c, l, r)$ and $(\cdot, 1, a', c', l', r')$.
%         \item Distributors $d : x(y + z) \simeq xy + xz$ and $d' : (x + y)z \simeq xz + yz$.
%         \item Absorbers $n : x0 \simeq 0$ and $n' : 0x \simeq 0$.
%     \end{enumerate}
%     Satisfying a number of coherence conditions
% \end{definition}

% We will call such a rig category a ring category if the monoidal structure $(+, 0, a, c, l, r)$ defines a $2$-group.

% Let $M$ be a symmetric monoidal category.
% \begin{definition}
%     $M$ is \textbf{traced} if there is a family of functions 
%     \[ (\_ \ominus U ): \Hom(X \oplus U, Y \oplus U) \to \Hom(X, Y)\]
%     satisfying:
%     \[ (f \circ (g \oplus U)) \ominus U = (f \ominus U) \circ g \]
%     \[ (g \oplus U) \circ f \ominus U = g \circ (f \ominus U) \]
%     For any $f : X \oplus U \to Y \oplus V$  :
%     \[ (Y \oplus g) \circ f \ominus U = f \circ (X \oplus g) \ominus V\] 
%     For any $f : X \oplus I \to Y \oplus I$ :
%     \[f \ominus 0 = \rho_Y \circ f \circ \rho_X^{-1}\]
%     For any $f : X \oplus U \oplus V \to Y \oplus U \oplus V$ 
%     \[ (f \ominus V) \ominus U = f \ominus (U \oplus V)\]
%     For any $f : X \oplus U \to Y \oplus U$ and $g : W \to Z$
%     \[ g \oplus (f \ominus U) = (g \oplus f) \ominus U\]
%     \[ \mathrm{sym}_{X} \ominus X = \id_X \]
% \end{definition}


% From this we can form the Int construction by defining a new category 
% \begin{definition}
%     $\mathrm{Int}(M)$ has objects given by pairs of objects in $M$ and
%     morphisms 
%     \[ f : (X^+, X^-) \to (Y^+, Y^-) \]
%     \[ f : X^+ \oplus Y^- \to Y^+ \oplus X^-  \]
%     Composition is defined using $\ominus$
% \end{definition}

% Let $W$ be the class of maps $f$ in $\mathrm{Int}(M)$ such that $f$ is an isomorphism
% in $M$.
% We consider the localisation of $M$ at $W$.

% Note $W$ is a very good class: It is a left and right calculus of fractions.

% Call this category $G(M)$.

% \begin{definition}
%     $G(M)$ has a symmetric monoidal product structure:
%     \[ \oplus : G(M) \times G(M) \to G(M) \]
%     \[ (X^+, X^-) \oplus (Y^+, Y^-) = (X^+ \oplus Y^+, X^- \oplus Y^-) \]
% \end{definition}

% Furthermore if $M$ is actually a rig category then $G(M)$ has a second symmetric monoidal structure.
% \[ \otimes : G(M) \times G(M) \to G(M) \]
% \[ X \otimes Y := (X^+ \otimes Y^+ \oplus X^- \otimes Y^-, X^+ \otimes Y^- \oplus X^- \otimes Y^+)\]
% With unit given by $(I, 0)$.


\section{The Category of Virtual Sets} 


\begin{definition}
    Given an injective map $f : X + A \to Y + A$, where $A$ is finite, we define a new map 
    \[ (f - A) : X \to Y \] 
    \[  (f-A)(x) = \begin{cases}
        y & \text{ if } f(x) = y : Y \\
        f_A(a) & \text{ if } f(x) = a : A 
    \end{cases} \]

    Where we define a helper for each $S \subseteq A$ 
    \[ f_S : S \to Y \]
    \[ f_S(a) = \begin{cases}
        y & \text{ if } f(a) = y : Y \\
        f_{S \backslash a}(a')&  \text{ if } f(a) = a' : A
    \end{cases}\]

    Note by injectivity of $f$ we have that $a$ and $a'$ are different, 
    and since $A$, we will eventually hit an element of $Y$.
\end{definition}

\begin{theorem}
  The operation $(\_ - A)$ forms a trace on the category of finite sites with injections.

  That is:
  \begin{enumerate}
    \item For any $f : Y + A \to Z + A$ and $g : X \to Y$ we have: 
    \[ (f \circ (g + A)) - A = (f - A) \circ g \]
    \item For any $f : Y \to Z$ and $g : X + A \to Y + A$ we have: 
    \[ ((f + A) \circ g) - A = f \circ (g - A) \]
    \item 
    For any $f : X + A \to Y + B$ and $g : B \to A$
    \[ ((Y + g) \circ f) - A = (f \circ (X + g)) - B\]
    \item
    For any $f : X + 0 \to Y \oplus 0$:
    \[f - 0 = \rho_Y \circ f \circ \rho_X^{-1}\]
    where $\rho_Z : Z \to Z + 0$ is the natural isomorphism.
    \item 
    For any $f : X + A + B \to Y + A + B$:
    \[ (f - B) - A = f - (A + B)\]
    \item
    For any $f : X + A \to Y + A$ and $g : W \to Z$
    \[ g + (f - A) = (g + f) - A\]
  \end{enumerate}
\end{theorem}
\begin{proof}
  We present a partial proof in the following lemmas.
\end{proof}

\begin{lemma}
    Given $f : X \to Y$ we have $(f + A) - A = f$.
\end{lemma}
\begin{proof}
    For any $x : X$ we have $(f + A)(\inl x) = \inl f(x) : Y$.
    Hence we see by definition that  $((f + A) - A)(x) = f(x)$.
\end{proof}

\begin{lemma}
    Given $f : X + A \to Y + A$ with $A$ finite, for any set $B$ we have
    \[ (f - A) + B = (f + B) - A \]
\end{lemma}
\begin{proof}
    Note for any $b : B$ we have $((f + B) - A)(b) = b$ by definition.
    and for $x : X$ we have $(f + B)(x) = f(x)$, and similarly for any $a : A$. 
    Unwinding definitions makes it clear that $f_A$ and $(f + B)_A$ are the same function,
    and further that $((f + B) - A)(x) = (f - A)(x)$.
\end{proof}

\begin{lemma}
    Given $f : X + A + B \to Y + A + B$ we have $(f - B) - A = f - (A + B)$.
\end{lemma}
\begin{proof}
    Need to come up with a better way of seeing this but we can redefine the notion of subtraction in the following way:
    For an easier way to see this think of a function $X + A \to Y + A$ as a kind of graph on the vertex set $X + A + Y + B$, with edges $x \to f(x)$.
    Then we can take a graph quotient identifying the two copies of $A$ and the left copy of $A$ with its image under $f$
    So $\inl a \sim \inr a$ and $\inl a ~ f(a)$.
    Taking thsi graph quotient results in the graph of the function $f - A : X \to Y$. 
    This is a kind of ``small step'' reduction process to produce $f - A$.
    Then we can see clearly that $f - B - A = f - (A + B)$ by just checking the two equivalence relations are the same, which they clearly are.
\end{proof}

\begin{lemma}
    We have $(g + A) \circ f - A = g \circ (f - A)$.
    Similarly $g \circ (f + A) = (g - A) \circ f$.
\end{lemma}


Given these we are now in a position to work with virtual sets.

\begin{definition}
    Define a category by 
    \begin{itemize}
        \item Objects pairs of finite sets $(X^+, X^-)$
        \item Morphisms $X \to Y$ given by an injective map $X^+ + Y^- \to Y^+ + X^-$.
        \item Identities given by identity maps
        \item Composition given by the following construction. Given $f : X^+ + Y^- \to Y^+ + X^-$ and $g : Y^+ + Z^- \to Z^+ + Y^-$ we define a map
        \begin{align*}
            h : X^+ + Z^- + Y^- &\xrightarrow{f + Z^-} Y^+ + Z^- + X^- \\ 
                            &\xrightarrow{g + X^-} Z^+ + Y^- + X^- \\
                            &\xrightarrow{\; \; \; \; \sim \; \; \; \;  }  Z^+ + X^- + Y^- 
        \end{align*}
        We will continue to abuse coherences and just write $h$ as $(g + X^-) \circ (f + Z^-)$.       
        The composition is given by $g \circ f := h - Y^-$. 
    \end{itemize}
\end{definition}

We now check that composition is associative.
\begin{proposition}
    Composition is associative.
\end{proposition}
\begin{proof}
    Let $f : A \to B$, $g : B \to C$ and $h : C \to D$.
    Then we calculate using the lemmas before
    \[ g \circ f = (g + A^-) \circ (f + C^-) - B^-\]
    \begin{align*}
        h \circ (g \circ f) 
        &= (h + A^-) \circ (g \circ f + D^-) - C^- 
        \\&= (h + A^-) \circ ( ((g + A^-) \circ (f + C^-) - B^- ) + D^-) - C^-
        \\&= (h + A^-) \circ ( ((g + A^-) \circ (f + C^-) + D^- ) - B^-) - C^-
        \\&= (h + A^-) \circ  ((g + A^- + D^-) \circ (f + C^- + D^-) - B^-) - C^-
        \\&= (h + A^- + B^-) \circ ((g + A^- + D^-) \circ (f + C^- + D^-)) - B^- - C^-
        \\&= (h + A^- + B^-) \circ (g + A^- + D^-) \circ (f + C^- + D^-) - (C^- + B^-)
    \end{align*} 
    Which is precisely the natural ``triple composition'' given between three morphisms.

    Similarly we have 
    \[ h \circ g = (h + B^-) \circ (g + D^-) - C^- \]
    \begin{align*}
        (h \circ g) \circ f 
        &= (h \circ g + A^-) \circ (f + D^-) - B^- 
        \\&= (((h + B^-) \circ (g + D^-) - C^-) + A^-) \circ (f + D^-) - B^-
        \\&=  (((h + B^-) \circ (g + D^-) + A^-) - C^-) \circ (f + D^-) - B^-
        \\&= ((h + B^- + A^-) \circ (g + D^- + A^-) - C^-) \circ (f + D^-) - B^-
        \\&= ((h + B^- + A^-) \circ (g + D^- + A^-)) \circ (f + D^- + C^-) - C^- - B^-
        \\&=  (h + B^- + A^-) \circ (g + D^- + A^-) \circ (f + D^- + C^-) - (B^- + C^-)
        \\&=  (h + B^- + A^-) \circ (g + D^- + A^-) \circ (f + D^- + C^-) - (C^- + B^-)
    \end{align*} 
\end{proof}

\begin{remark}
    Here is a relatively convincing reason why this category might not be correct: two objects are only isomorphic if their pairs of objects are isomorphic.
    This is unlike the intuition that we would expect $(3, 2) \simeq (2, 1)$.

    We prove this now. Suppose $X \simeq Y$ in the category. 
    Suppose they don't have the same size (this is fine because of decidabilty).
    WLoG suppose $| Y^+ | < | X^+ |$.
    Call $f : X \to Y$ and $g : Y \to X$ the functions in the isomorphism.
    Then note that for some $x^+ : X^+$ we must have $f(x^+) \in X^-$, due to size.
    Then by definition $(g + X^-)( (f + X^-)(x^+) ) = f(x^+)$, since $g$ acts constantly on this copy of $X^-$. 
    Then clearly $(g \circ f)(x^+) = f(x^+) \neq x^+$. 
    So we must have $| Y^+ | \geq | X^+ |$
    Arguing dually we see that $|Y^+| = |X^+|$ and so we deduce $|Y^-| = |X^-|$.
\end{remark}

Composition behaves differently then we would expect on most morphisms.
But on morphisms given by a pair of injective maps everything works out correctly.

\begin{definition}
    Specifically given $f : X^+ \to Y^+$ and $g : Y^- \to X^-$ we obtain a map
    \[\langle f, g \rangle := f + g : X^+ + Y^- \to Y^+ + X^- = X \to Y \]
\end{definition}

\begin{lemma}
    Given $\langle f_0, g_0 \rangle : X \to Y$ and $\langle f_1, g_1 \rangle : Y \to Z$ we have
    \[\langle f_1, g_1 \rangle \circ \langle f_0, g_0 \rangle  = \langle f_1 \circ f_0, g_1 \circ g_0 \rangle\] 
\end{lemma}
\begin{proof}
    This is clear since $(f_1 +g_1 +X^-) \circ (f_0 + g_0 + Z^-)= (f_1 \circ f_0) + g_1 + g_0 : X^+ + Z^- + Y^- \to Z^+ + Y^- + X^-$
    Thus subtracting off ... \todo{}
\end{proof}

Let $W$ be the class of morphisms which are bijections as set functions.

\begin{definition}
  A \textbf{congruence} on a category $\cC$ is an equivalence relation $\sim_{a,b}$ on each $\hom_\cC(a, b)$ such that if $f_1 ~_{a, b} f_2$ and $g_1 ~_{b, c} g_2$ then $g_1 \circ f_1 ~_{a, c} g_2 \circ f_2$.
\end{definition}

Given a family of relations $R_{a, b}$ on $\hom_\cC(a, b)$ we can create the congruence generated by $R$ by taking the symmetric, reflexive, transitive and composition closure of $R$.

Let $~$ be the congreunce relation on $\cC$ generated by for all $f \in W$ that $f \circ f^{-1} \sim \id$. 

We define the category $\VSet := \cC / \sim$ to be the quotient of $\cC$ by this congreunce relation.

\begin{definition}
    The category $\cC$ has two monoidal products on it, given by addition and multiplication.
    Specifically
    \[ \oplus : \cC \times \cC \to \cC \]
    \[ X \oplus Y := (X^+ + Y^+, X^- + Y^-) \]
    Functorial action is given by identity on the added component
    The associator is given by the isomorphism on $\cC$ by the pair of associators on each component.
    The monoidal unit is given by $(0, 0) \in \cC$, with similar unitors given by pairs of unitors.

    Since composition of pairs acts just as composition pairwise, the coherence conditions reduce to coherence conditions in $\Set$.
    
    The second monoidal product is given by 
    \[ X \otimes Y := (X^+ \times Y^+ + X^- \times Y^-, X^+ \times Y^- + X^- \times Y^+)\]
    With unit given by $(1, 0)$.
\end{definition}

\begin{theorem}
  These monoidal products descend to the category $\VSet$ with the quotient $\cC \to \VSet$ being a monoidal functor.
\end{theorem}

\end{document}